%%%%%%%%%%%%%%%%%%%%%%%%%%%%%%%%%%%%%%%%%%%%%%%%%%
% PREAMBLE
%%%%%%%%%%%%%%%%%%%%%%%%%%%%%%%%%%%%%%%%%%%%%%%%%%


%%%%%%%%%%%%%%%%%%%%%%%%%%%%%%%%%%%%%%%%%%%%%%%%%%
% SPRACHE UND SCHRIFT
%%%%%%%%%%%%%%%%%%%%%%%%%%%%%%%%%%%%%%%%%%%%%%%%%%

% UTF-8-Zeichenkodierung der Quelldateien
\usepackage[utf8]{inputenc}

% Deutsche Silbentrennung, Übersetzungen und Typografie
\usepackage[ngerman]{babel}

% T1-Schriftkodierung für korrekte Darstellung europäischer Zeichen
\usepackage[T1]{fontenc}

% Latin-Modern-Schrift
\usepackage{lmodern}
% Sprachabhängige Anführungszeichen
\usepackage{csquotes}

%%%%%%%%%%%%%%%%%%%%%%%%%%%%%%%%%%%%%%%%%%%%%%%%%%
% SEITENLAYOUT UND ABSTÄNDE
%%%%%%%%%%%%%%%%%%%%%%%%%%%%%%%%%%%%%%%%%%%%%%%%%%

% Einstellung von Seitenrändern und Papierformat
\usepackage{geometry}
\geometry{
    a4paper,
    margin=3cm
}

% Ermöglicht die Einstellung des Zeilenabstands
\usepackage[onehalfspacing]{setspace}

% Verbesserte Positionierung von Floats
\usepackage{dblfloatfix}

% Ermöglicht Bilder, die von Text umflossen werden
\usepackage{wrapfig}

% Blocksatz und zusätzliche Textausrichtungen
\usepackage{ragged2e}

%%%%%%%%%%%%%%%%%%%%%%%%%%%%%%%%%%%%%%%%%%%%%%%%%%
% TABELLEN
%%%%%%%%%%%%%%%%%%%%%%%%%%%%%%%%%%%%%%%%%%%%%%%%%%

% Tabellen mit automatisch angepasster Spaltenbreite
\usepackage{tabularx}

% Tabellen über mehrere Seiten
\usepackage{longtable}

% Erweiterte Tabellenfunktionen und eigene Spaltentypen
\usepackage{array}

% Tabellenzellen über mehrere Zeilen
\usepackage{multirow}


%%%%%%%%%%%%%%%%%%%%%%%%%%%%%%%%%%%%%%%%%%%%%%%%%%
% LISTEN UND AUFZÄHLUNGEN
%%%%%%%%%%%%%%%%%%%%%%%%%%%%%%%%%%%%%%%%%%%%%%%%%%

% Erweiterte Nummerierung von enumerate-Umgebungen
\usepackage{enumerate}

% Flexible Konfiguration von itemize, enumerate und description
\usepackage{enumitem}


%%%%%%%%%%%%%%%%%%%%%%%%%%%%%%%%%%%%%%%%%%%%%%%%%%
% MATHEMATIK UND THEOREME
%%%%%%%%%%%%%%%%%%%%%%%%%%%%%%%%%%%%%%%%%%%%%%%%%%

% Erweiterte mathematische Umgebungen und Befehle
\usepackage{amsmath}

% Definition und Nummerierung von Sätzen, Definitionen etc.
\usepackage[
    thmmarks,
    amsmath,
    hyperref,
    noconfig
]{ntheorem}


%%%%%%%%%%%%%%%%%%%%%%%%%%%%%%%%%%%%%%%%%%%%%%%%%%
% BILDER UND GRAFIKEN
%%%%%%%%%%%%%%%%%%%%%%%%%%%%%%%%%%%%%%%%%%%%%%%%%%

% Einfügen und Skalieren von Bildern
\usepackage{graphicx}


%%%%%%%%%%%%%%%%%%%%%%%%%%%%%%%%%%%%%%%%%%%%%%%%%%
% FARBEN
%%%%%%%%%%%%%%%%%%%%%%%%%%%%%%%%%%%%%%%%%%%%%%%%%%

% Erweiterte Farbunterstützung, auch für Tabellen
\usepackage[table]{xcolor}

% Eigene Farben
\definecolor{codepurple}{rgb}{0.58,0,0.82}

% Telekom-Magenta
% RGB-Werte müssen zwischen 0 und 1 liegen.
\definecolor{magenta}{RGB}{226,0,116}

% Hintergrundfarbe für Code-Listings
\definecolor{backcolour}{RGB}{245,245,245}


%%%%%%%%%%%%%%%%%%%%%%%%%%%%%%%%%%%%%%%%%%%%%%%%%%
% QUELLCODE
%%%%%%%%%%%%%%%%%%%%%%%%%%%%%%%%%%%%%%%%%%%%%%%%%%

% Quellcode mit Syntax-Highlighting über Pygments
% Benötigt bei pdflatex normalerweise -shell-escape
\usepackage{minted}

% Quellcode direkt mit LaTeX darstellen
\usepackage{listings}

% Globale Einstellungen für listings
\lstset{
    basicstyle=\ttfamily\footnotesize,
    firstnumber=1,
    stepnumber=1,
    numbersep=8pt,
    numbers=left,
    breaklines=true,
    frame=lines,
    captionpos=b,
    numberstyle=\tiny\color{gray},
    columns=fullflexible,
    showstringspaces=false,
    backgroundcolor=\color{backcolour},
}


%%%%%%%%%%%%%%%%%%%%%%%%%%%%%%%%%%%%%%%%%%%%%%%%%%
% LITERATURVERZEICHNIS
%%%%%%%%%%%%%%%%%%%%%%%%%%%%%%%%%%%%%%%%%%%%%%%%%%

% Literaturverwaltung mit Biber im IEEE-Stil
\usepackage[
    backend=biber,
    style=ieee,
    sorting=none
]{biblatex}

% BibTeX-Datei
\addbibresource{./quellen/quellen.bib}


% Formatierung des BibTeX-Schriftzugs
\def\BibTeX{%
    {\rm B\kern-.05em%
            {\sc i\kern-.025em b}\kern-.08em%
        T\kern-.1667em%
        \lower.7ex\hbox{E}\kern-.125emX}%
}


%%%%%%%%%%%%%%%%%%%%%%%%%%%%%%%%%%%%%%%%%%%%%%%%%%
% DOKUMENTSTRUKTUR
%%%%%%%%%%%%%%%%%%%%%%%%%%%%%%%%%%%%%%%%%%%%%%%%%%

% Einbinden einzelner LaTeX-Dateien als Subfiles
\usepackage{subfiles}

% Verwaltung eines Abkürzungsverzeichnisses
% printonlyused zeigt nur tatsächlich verwendete Abkürzungen
\usepackage[printonlyused]{acronym}


%%%%%%%%%%%%%%%%%%%%%%%%%%%%%%%%%%%%%%%%%%%%%%%%%%
% URLs UND HYPERLINKS
%%%%%%%%%%%%%%%%%%%%%%%%%%%%%%%%%%%%%%%%%%%%%%%%%%

% Verbesserte Darstellung und Zeilenumbrüche von URLs
\usepackage{url}

% Klickbare Links, Verweise und Literaturzitate
% Sollte möglichst spät geladen werden.
\usepackage[
    colorlinks=true,
    urlcolor=blue,
    linkcolor=blue,
    citecolor=codepurple,
    menucolor=red,
    runcolor=cyan
]{hyperref}


%%%%%%%%%%%%%%%%%%%%%%%%%%%%%%%%%%%%%%%%%%%%%%%%%%
% THEOREME
%%%%%%%%%%%%%%%%%%%%%%%%%%%%%%%%%%%%%%%%%%%%%%%%%%

% Satz-Umgebung mit Nummerierung nach Abschnitt
% Beispiel: Satz 2.1
\newtheorem{satz}{Satz}[section]

% Normaler statt kursiver Text innerhalb der Theorem-Umgebungen
\theorembodyfont{\upshape}


%%%%%%%%%%%%%%%%%%%%%%%%%%%%%%%%%%%%%%%%%%%%%%%%%%
% EIGENE TABELLENTYPEN
%%%%%%%%%%%%%%%%%%%%%%%%%%%%%%%%%%%%%%%%%%%%%%%%%%

% Eigener Spaltentyp "x"
\newcolumntype{x}[1]{%
    !{\centering\arraybackslash\vrule width #1}%
}


%%%%%%%%%%%%%%%%%%%%%%%%%%%%%%%%%%%%%%%%%%%%%%%%%%
% KOPF- UND FUSSZEILE
%%%%%%%%%%%%%%%%%%%%%%%%%%%%%%%%%%%%%%%%%%%%%%%%%%

\usepackage{scrlayer-scrpage}

% Bestehende Kopf- und Fußzeilen entfernen
\clearpairofpagestyles

% Zentrierte Seitenzahl
\cfoot*{\pagemark}

% Seitenstil aktivieren
\pagestyle{scrheadings}